\documentclass[12pt,a4paper]{article}
\usepackage[utf8]{inputenc}
\usepackage[russian]{babel}
\usepackage{amsmath, amssymb, amsthm}
\usepackage{graphicx}
\usepackage{hyperref}
\usepackage{listings}
\usepackage{xcolor}
\usepackage{geometry}
\usepackage{float}
\usepackage{caption}
\usepackage{subcaption}
\usepackage{enumitem}
\usepackage{titlesec}
\usepackage{fancyhdr}
\usepackage{array}
\usepackage{longtable}
\usepackage{booktabs}
\usepackage{ragged2e}

% Настройка полей
\geometry{
    a4paper,
    left=2.5cm,
    right=1.5cm,
    top=2cm,
    bottom=2cm
}

% Настройка заголовков
\titleformat{\section}{\Large\bfseries}{\thesection}{1em}{}
\titleformat{\subsection}{\large\bfseries}{\thesubsection}{1em}{}
\titleformat{\subsubsection}{\normalsize\bfseries}{\thesubsubsection}{1em}{}

% Настройка кода
\lstset{
    language=Python,
    basicstyle=\ttfamily\small,
    keywordstyle=\color{blue},
    commentstyle=\color{green!40!black},
    stringstyle=\color{red},
    numbers=left,
    numberstyle=\tiny,
    breaklines=true,
    frame=single,
    captionpos=b,
    showstringspaces=false,
    tabsize=4
}

% Настройка колонтитулов
\pagestyle{fancy}
\fancyhf{}
\fancyhead[L]{Лабораторная работа №2: Развитие AI-сервиса}
\fancyhead[R]{\thepage}
\fancyfoot[C]{}

\begin{document}

\begin{titlepage}
    \begin{center}
        \textbf{МИНИСТЕРСТВО НАУКИ И ВЫСШЕГО ОБРАЗОВАНИЯ РОССИЙСКОЙ ФЕДЕРАЦИИ}
        \vspace{0.5cm}

        \textbf{Федерального государственного автономного образовательного учреждения \\ высшего образования}
        \vspace{0.3cm}

        \textbf{«НАЦИОНАЛЬНЫЙ ИССЛЕДОВАТЕЛЬСКИЙ ТЕХНОЛОГИЧЕСКИЙ УНИВЕРСИТЕТ МИСИС»}
        \vspace{2cm}


        \textbf{Институт компьютерных наук}
        \vspace{0.3cm}

        \textbf{Кафедра Инженерной кибернетики}
        \vspace{1cm}

        \textbf{Лабораторная работа №2}
        \vspace{0.3cm}

        \textbf{по дисциплине «Современные интеллектуальные сетевые сервисы»}
        \vspace{1.5cm}

        \huge{\textbf{Развитие AI-сервиса: \\ сквозная идентификация запросов, \\ асинхронная обработка и интеграция LLM}}
        \vspace{2cm}
    \end{center}

    \begin{flushright}
        \textbf{Выполнил:} \\
        Студент группы \\
        И.О. Фамилия \\
        \vspace{0.5cm}
        \textbf{Проверил:} \\
        Преподаватель \\
        И.О. Фамилия \\
        \vspace{2cm}
        \textbf{Оценка:} \hrulefill
    \end{flushright}

    \begin{center}
        \vfill
        Москва, 2026
    \end{center}
\end{titlepage}

\tableofcontents
\newpage

\section{Цель работы}

Цель практической работы — развить ранее созданный REST API сервиса агроскоринга и приблизить его архитектуру к реальному промышленному AI-сервису.

В ходе работы необходимо изучить и реализовать механизмы, которые возникают при эксплуатации сервиса: сквозную идентификацию запросов с помощью \texttt{request\_id}, использование промежуточных слоёв (middleware), управление очередью запросов и пакетным инференсом, хранение результатов в базе данных SQLite, асинхронное взаимодействие с внешними сервисами и подключение большой языковой модели через HTTP API.

Последовательный поток обработки запроса представлен на рисунке~\ref{fig:flow}.

\begin{figure}[H]
    \centering
    \includegraphics[width=0.45\textwidth, keepaspectratio]{request_flow.png}
\begin{figure}
        \centering
        \includegraphics[width=0.5\linewidth]{request_flow.png}
        \caption{Enter Caption}
        \label{fig:placeholder}
    \end{figure}
        \caption{Обработка запроса в проектируемом API AI-сервиса агроскоринга}
    \label{fig:flow}
\end{figure}

\section{Задачи работы}

Стабильность API агроскоринга, спроектированного в лабораторной работе №1, обеспечивалась только контрактом (Pydantic-схемами данных). Теперь необходимо разработать набор механизмов, которые позволяют сервису работать под нагрузкой, переживать сбои и быть наблюдаемым.

В ходе лабораторной работы №2 необходимо:

\begin{enumerate}
    \item Модифицировать существующий API агроскоринга таким образом, чтобы каждому HTTP-запросу назначался уникальный \texttt{request\_id}.
    \item Реализовать промежуточный слой ПО middleware, который:
    \begin{itemize}
        \item принимает \texttt{X-Request-ID} от клиента;
        \item при его отсутствии генерирует UUID (Universally Unique Identifier);
        \item сохраняет идентификатор в контексте запроса;
        \item возвращает его клиенту в HTTP-заголовке;
        \item измеряет время обработки запроса.
    \end{itemize}
    \item Реализовать два режима инференса: одиночный и пакетный.
    \item Организовать очередь запросов с использованием механизма \texttt{asyncio.Queue}.
    \item Реализовать динамическое формирование пакета запросов (dynamic batching) — автоматическое объединение нескольких одиночных HTTP-запросов в один пакетный вызов модели.
    \item Задать эксплуатационные ограничения: максимальный размер batch, максимальное время ожидания формирования batch, максимальный размер очереди, таймаут инференса.
    \item Реализовать механизм управления потоком запросов (backpressure), при котором система, не успевающая обрабатывать входящие запросы, сигнализирует клиенту, что нужно снизить скорость отправки или временно прекратить отправку. В данном проекте достаточно вернуть клиенту корректный HTTP-код при переполнении очереди.
    \item Подключить базу данных SQLite для постоянного хранения результатов прогнозирования.
    \item Для асинхронной работы с SQLite использовать драйвер \texttt{aiosqlite}.
    \item Реализовать: запись одиночного результата, пакетную запись результатов, чтение результатов, поиск по \texttt{request\_id}, пагинацию результатов.
    \item Подключить внешнюю LLM через HTTP API.
    \item Для асинхронного HTTP-вызова использовать \texttt{httpx.AsyncClient}.
    \item Реализовать endpoint для непосредственного запроса к LLM.
    \item Реализовать обработку ошибок внешних зависимостей: ошибка БД, таймаут модели, переполнение очереди, недоступность LLM, ошибка внешнего LLM API.
    \item Протестировать разработанный сервис через Swagger UI и проанализировать полный путь одного запроса от клиента до модели, БД и внешней LLM.
\end{enumerate}

\section{Теоретическая часть}

\subsection{Что такое request\_id}

Если приложение небольшое, можно посмотреть на лог и примерно понять, какой запрос выполнялся. В нагруженной системе одновременно могут выполняться сотни запросов, их сообщения в журнале будут перемешаны. Поэтому каждому входящему HTTP-запросу назначается уникальный идентификатор: \texttt{request\_id}. Он используется для отслеживания конкретного HTTP-запроса на протяжении всего его пути через систему.

Существуют два варианта создания \texttt{request\_id}:

\begin{enumerate}
    \item Клиент сам может передать \texttt{X-Request-ID} — вспомогательный опциональный заголовок HTTP, который содержит уникальный идентификатор запроса. Это удобно, если запрос проходит через несколько систем, например CRM, API Gateway, Database. Все системы могут использовать один идентификатор.
    \item Если клиент его не передал, сервер создаёт свой:
\end{enumerate}

\begin{lstlisting}[caption={Генерация request\_id}]
request_id = str(uuid.uuid4())
\end{lstlisting}

\subsection{Middleware}

Чтобы не писать код генерации \texttt{request\_id} в каждом endpoint, можно использовать middleware.

Middleware в FastAPI — это промежуточный слой (функция/класс), через который проходит каждый HTTP-запрос перед тем, как попасть в обработчик endpoint, и каждый ответ перед тем, как уйти клиенту. Это даёт единую точку для обработки ошибок, логирования и проверки аутентификации, а также позволяет гибко добавлять и менять компоненты без правок в ядре приложения.

Middleware выполняется:

\begin{enumerate}
    \item до endpoint;
    \item передаёт управление endpoint;
    \item получает сформированный response;
    \item может изменить response;
    \item возвращает его клиенту.
\end{enumerate}

Например:

\begin{lstlisting}[caption={Заготовка middleware}]
@app.middleware("http")
async def request_context_middleware(
    request: Request,
    call_next
):
    ...
\end{lstlisting}

Здесь \texttt{request} — входящий HTTP-запрос, а \texttt{call\_next(request)} — передача запроса следующему компоненту приложения.

\subsubsection{Хранение данных в request.state}

FastAPI позволяет сохранить информацию, относящуюся только к текущему запросу, для этого используют \texttt{request.state}. В него можно записывать произвольные данные, а затем читать их в разных частях приложения во время обработки одного и того же запроса.

Обычно \texttt{request\_id} генерируется в middleware. Если в заголовке запроса уже есть \texttt{X-Request-ID}, используется оно. Если нет — генерируется новый UUID. Затем значение сохраняется в \texttt{request.state.request\_id}.

\subsubsection{Измерение времени запроса}

Middleware удобно использовать для измерения времени обработки:

\begin{equation}
T_{\text{request}} = t_{\text{response}} - t_{\text{request}}
\end{equation}

Например:

\begin{lstlisting}[caption={Измерение времени обработки}]
start_time = time.perf_counter()
response = await call_next(request)
process_time = time.perf_counter() - start_time
\end{lstlisting}

Получаем простейшую метрику latency.

\subsection{Что такое batching}

Батчинг (batching, или пакетная обработка) — метод повышения производительности за счёт сокращения количества отдельных вызовов модели, который представляет собой группирование нескольких похожих входных запросов в один общий вызов инференса. Существует два типа методов батчинга: статический и динамический.

Batching позволяет уменьшить относительную стоимость некоторых фиксированных операций, например передачи данных, постпроцессинга, обмена данными между CPU и GPU. Поэтому batching часто увеличивает пропускную способность (throughput) — число обработанных объектов в единицу времени.

\subsubsection{Explicit batching}

Необходимо различать явный (explicit) и динамический (dynamic) batching.

Explicit batching — это подход, при котором система явно управляет пакетной обработкой, то есть сама формирует пакеты запросов или операций и инициирует их выполнение.

Сам клиент передаёт несколько объектов:

\begin{lstlisting}[caption={Explicit batch-запрос}, language=Python]
POST /batch-predict

{
    "items": [
        {"farm_id": "FARM-001"},
        {"farm_id": "FARM-002"},
        {"farm_id": "FARM-003"}
    ]
}
\end{lstlisting}

Сервис сразу знает состав batch.

\subsubsection{Dynamic batching}

Dynamic batching, напротив, — это подход, при котором система автоматически формирует пакеты запросов на лету, адаптируясь к текущей нагрузке и условиям.

Клиенты отправляют обычные одиночные запросы, они накапливаются в очереди, и когда система видит возможность сформировать пакет (достигнут определённый размер или прошло заданное время), она отправляет его на обработку.

Таким образом, для клиента API запрос остаётся одиночным, а сервер самостоятельно оптимизирует вычисления.

Dynamic batching оптимален там, где нагрузка неравномерна, а задержка ответа критична: онлайн-инференс, API-сервисы, чат-боты, LLM-серверы.

\subsubsection{Очередь запросов}

Для dynamic batching используется \texttt{asyncio.Queue}.

\texttt{asyncio.Queue} — это асинхронная FIFO («первым пришёл — первым ушёл») очередь. Основные методы и особенности работы асинхронной очереди можно посмотреть в документации: \url{https://docs.python.org/3/library/asyncio-queue.html}.

Для работы с очередью нужна специальная функция. Воркер (worker) — это долгоживущая корутина или процесс, который извлекает задачи из очереди и выполняет их. У worker много задач, но основная — пакетная обработка.

\subsubsection{Формирование dynamic batch}

Для пакетной обработки зададим два параметра:

\begin{lstlisting}[caption={Параметры batch}]
MAX_BATCH_SIZE = 8
MAX_BATCH_WAIT_MS = 50
\end{lstlisting}

Worker начинает формировать пакет после первого запроса и запускает модель, если выполнено хотя бы одно условие:

\begin{equation}
\text{batch\_size} = \text{MAX\_BATCH\_SIZE}
\quad \text{ИЛИ} \quad
\text{время ожидания} \ge \text{MAX\_BATCH\_WAIT\_MS}
\end{equation}

Ожидание только наполнения пакета может значительно увеличить задержку (latency) при небольшой нагрузке.

\subsection{Latency и throughput}

При проектировании batching возникает компромисс между двумя показателями.

Задержка (latency) — время обработки одного запроса:

\begin{equation}
\text{Latency} = t_{\text{response}} - t_{\text{request}}
\end{equation}

Пропускная способность (throughput) — число обработанных запросов за единицу времени:

\begin{equation}
\text{Throughput} = \frac{N}{T}
\end{equation}

Больший пакет позволяет повысить пропускную способность, но может увеличить задержку, поэтому достичь компромисса между этими двумя метриками — одна из основных инженерных задач при эксплуатации AI-сервисов.

На практике этот компромисс управляется несколькими рычагами:

\begin{itemize}
    \item Максимальный размер батча (\texttt{MAX\_BATCH\_SIZE}) — верхняя граница, ограниченная объёмом памяти GPU и требованиями к латентности.
    \item Окно ожидания (\texttt{MAX\_BATCH\_WAIT\_MS}/timeout) — сколько времени сервис готов ждать накопления батча, прежде чем отправить его на выполнение. Типичные значения — единицы или десятки миллисекунд.
    \item Целевая латентность (SLO) — например, p95 или p99 времени ответа, под которую подстраивается размер батча.
    \item Приоритизация запросов — возможность обрабатывать часть запросов вне очереди, если для них критична задержка.
    \item Адаптивные стратегии — размер батча и окно ожидания динамически меняются в зависимости от текущей нагрузки, чтобы не перегружать GPU в пик и не простаивать в спад.
\end{itemize}

Итоговая задача инженера — найти такую конфигурацию, при которой сервис укладывается в заданные SLO по задержке и одновременно выжимает максимум из доступного железа.

\subsection{Механизм управления потоком запросов backpressure}

Предположим, модель обрабатывает 100 запросов/с, а клиенты отправляют 1000 запросов/с.

Если очередь ничем не ограничена, тысячи задач могут накапливаться в памяти, это приводит к росту RAM, росту latency, таймаутам.

Нужен механизм, ограничивающий очередь, например:

\begin{lstlisting}[caption={Ограничение очереди}]
QUEUE_MAX_SIZE = 100
\end{lstlisting}

Если очередь заполнена, новый запрос отклоняется.

Система сообщает потребителю: «Сейчас я не могу принимать задачи быстрее, чем успеваю их обрабатывать».

\subsection{FastAPI и асинхронность}

FastAPI представляет собой решение для разработки API, опирающееся на асинхронную модель обработки запросов.

Большинство AI-сервисов являются I/O-bound, они ждут ответов от внешних сервисов, поэтому библиотека \texttt{asyncio} — естественный выбор для таких сервисов.

В основе \texttt{asyncio} — event loop (цикл событий). Это бесконечный цикл, который:

\begin{enumerate}
    \item Берёт задачу из очереди.
    \item Выполняет её до первого \texttt{await}.
    \item Отправляет задачу ждать результата.
    \item Берёт следующую задачу из очереди.
    \item Возвращается к ожидающим задачам, как только они получили данные.
\end{enumerate}

Корутина (coroutine) — это асинхронная функция, объявленная с \texttt{async def}. При вызове она возвращает объект корутины, который нужно запустить через событийный цикл. Это позволяет эффективно обслуживать множество одновременно ожидающих запросов (БД, сеть, файл, HTTP).

\texttt{await} — точка ожидания, она сообщает событийному циклу: «дождись здесь результата, а пока переключись на другие задачи».

В нашем приложении эндпоинты определяются как корутины.

Одна из распространённых ошибок — пропуск \texttt{await} при вызове асинхронной функции. Python не сообщит об ошибке, а просто вернёт объект корутины вместо ожидаемого результата. Поэтому внутри async-функций всегда следует использовать \texttt{await}.

Существуют высокоуровневые функции \texttt{asyncio.gather()}, \texttt{asyncio.wait()}, которые используются для запуска нескольких задач параллельно.

\subsection{SQLite}

SQLite — встраиваемая реляционная база данных.

В отличие от PostgreSQL или MySQL для неё не требуется отдельный сервер. База хранится в файле \texttt{agro\_scoring.db}.

В таблице \texttt{predictions} можно хранить, например: \texttt{request\_id}, \texttt{farm\_id}, \texttt{risk\_score}, \texttt{risk\_level}, \texttt{recommendation}, \texttt{model\_version}.

Нельзя напрямую использовать СУБД \texttt{sqlite3} внутри \texttt{async def}, так как она имеет синхронный интерфейс.

Например, \texttt{cursor.execute()}, \texttt{conn.commit()} являются блокирующими операциями. Если выполнить их непосредственно внутри \texttt{async def endpoint()}, то event loop будет занят до завершения операции.

\texttt{aiosqlite} — это асинхронный драйвер для работы с базой данных SQLite в Python, построенный на основе стандартного модуля \texttt{sqlite3}.

Подключение:

\begin{lstlisting}[caption={Подключение к SQLite}]
async with aiosqlite.connect("agro_scoring.db") as db:
    ...
\end{lstlisting}

Выполнение SQL:

\begin{lstlisting}[caption={Выполнение SQL}]
await db.execute(...)
\end{lstlisting}

Фиксация транзакции:

\begin{lstlisting}[caption={Фиксация транзакции}]
await db.commit()
\end{lstlisting}

Чтение:

\begin{lstlisting}[caption={Чтение одной строки}]
row = await cursor.fetchone()
\end{lstlisting}

Использование \texttt{aiosqlite} делает программный интерфейс асинхронным, но не изменяет фундаментальные свойства самой SQLite. SQLite остаётся файловой БД и имеет ограничения по конкурентной записи.

Будем переключать БД в режим конкурентной записи WAL (Write-Ahead Logging — журналирование с предварительной записью). В этом режиме читатели и писатели не блокируют друг друга: читатели видят согласованный снимок данных на момент начала транзакции, а писатель может продолжать работу, не дожидаясь завершения чтения.

\subsubsection{Параметризованные SQL-запросы}

Данные пользователя не следует подставлять в SQL через f-строку:

\begin{lstlisting}[caption={Небезопасный SQL}]
f"""
SELECT *
FROM predictions
WHERE request_id = '{request_id}'
"""
\end{lstlisting}

Правильный вариант:

\begin{lstlisting}[caption={Параметризованный SQL}]
await db.execute(
    """
    SELECT *
    FROM predictions
    WHERE request_id = ?
    """,
    (request_id,)
)
\end{lstlisting}

Такой запрос называется параметризованным. Он безопаснее, правильно экранирует строки, снижает риск SQL-инъекций, корректнее работает с типами.

\subsubsection{Транзакция и commit}

SQL-команда \texttt{INSERT INTO} сама по себе ещё не означает гарантированное сохранение результата. После записи необходимо выполнить:

\begin{lstlisting}[caption={Фиксация одиночной записи}]
await db.commit()
\end{lstlisting}

Для пакетной обработки можно использовать:

\begin{lstlisting}[caption={Пакетная запись}]
await db.executemany(...)
await db.commit()
\end{lstlisting}

В этом случае несколько записей сохраняются одной транзакцией, это эффективнее, чем выполнять отдельный \texttt{commit()} после каждого элемента пакета.

\subsubsection{Пагинация}

Пагинация — это способ отдавать данные порциями, а выбор конкретного метода — компромисс между удобством навигации и производительностью на больших объёмах.

Мы будем использовать offset-пагинацию. Offset-пагинация (она же \texttt{LIMIT \dots\ OFFSET}) — самый простой и исторически первый способ постраничного вывода данных. Идея проста: чтобы получить N-ю страницу, мы пропускаем $(N-1) \cdot \text{page\_size}$ записей и берём следующие \texttt{page\_size}.

Чтобы продемонстрировать механизм пагинации в нашем приложении, мы будем возвращать одну страницу прогнозов для выбранного дня недели.

\subsection{Подключение внешней LLM}

После получения результата агроскоринга его можно дополнительно объяснить естественным языком.

Например:

\begin{lstlisting}[caption={Пример результата}]
risk_score = 0.9
risk_level = high
\end{lstlisting}

может быть преобразовано в: «Полученная оценка соответствует высокой категории риска. Результат требует дополнительной проверки оператором в соответствии с правилами банка».

Для этого наш FastAPI-сервис выступает уже не только как сервер, но одновременно как HTTP-клиент другого API.

\subsubsection{OpenAI-совместимый формат LLM API}

Будем использовать Vireonix API, не требует авторизации, предоставляет бесплатные токены модели \texttt{auto} по API:

\begin{lstlisting}[caption={Endpoint LLM}]
POST /v1/chat/completions
\end{lstlisting}

Запрос имеет типовую структуру:

\begin{lstlisting}[caption={Структура запроса к LLM}]
{
    "model": "auto",
    "messages": [
        {
            "role": "system",
            "content": "..."
        },
        {
            "role": "user",
            "content": "..."
        }
    ]
}
\end{lstlisting}

\subsubsection{httpx.AsyncClient}

AI-сервис — это прежде всего множество HTTP-вызовов: к LLM API, к инструментам поиска, к внешним сервисам. Их нужно организовывать асинхронно, с повторными попытками вызова и таймаутами. Поэтому вместо синхронного \texttt{requests} мы возьмём \texttt{httpx.AsyncClient} — асинхронный HTTP-клиент из библиотеки HTTPX.

Все методы клиента (\texttt{get}, \texttt{post}, \texttt{put}, \texttt{delete} и другие) являются асинхронными. Чтобы их вызвать, нужно использовать ключевое слово \texttt{await}. Например:

\begin{lstlisting}[caption={Асинхронный вызов}]
response = await client.post(...)
\end{lstlisting}

Создание нового \texttt{httpx.AsyncClient()} для каждого HTTP-запроса возможно, но не оптимально. Лучше создать один клиент при старте приложения и закрыть при остановке.

\subsection{Работа с ошибками внешних компонентов}

Когда API начинает зависеть от других систем, появляются новые категории ошибок.

\begin{longtable}{|p{2.5cm}|p{10.5cm}|}
\caption{HTTP-коды ошибок внешних компонентов}
\label{tab:errors}\\
\hline
\textbf{HTTP-код} & \textbf{Ситуация} \\
\hline
\endfirsthead
\hline
\textbf{HTTP-код} & \textbf{Ситуация} \\
\hline
\endhead
400 & нарушение бизнес-правил \\
404 & прогноз не найден \\
422 & ошибка Pydantic \\
429 & очередь переполнена \\
500 & необработанная внутренняя ошибка \\
502 & внешний LLM API вернул некорректный ответ \\
503 & БД, модель или внешний сервис недоступны \\
504 & превышен timeout инференса или LLM \\
\hline
\end{longtable}

Если наш FastAPI работает корректно, но внешний LLM-провайдер отвечает ошибкой, наш сервис выполняет роль посредника и 502 Bad Gateway лучше отражает причину сбоя.

Чтобы зависший внешний сервис не удерживал ресурсы приложения, необходимо установить таймаут:

\begin{lstlisting}[caption={Таймаут LLM}]
LLM_TIMEOUT_SECONDS = 60.0
\end{lstlisting}

Если ответ не получен вовремя, то обработчик выдаст ошибку 504 Gateway Timeout.

\subsection{Lifespan в FastAPI}

Lifespan — это механизм, который позволяет (особенно в асинхронной среде) централизованно инициализировать и закрывать подключения к базе данных, кэшу или внешним API.

\section{Задания}

\subsection{Задание 1. Реализуйте сквозной request\_id}

Модифицируйте приложение таким образом, чтобы каждый HTTP-запрос имел уникальный идентификатор.

Требования:

\begin{enumerate}
    \item Если клиент передал \texttt{X-Request-ID} — использовать его.
    \item Если \texttt{X-Request-ID} отсутствует — сгенерировать UUID.
    \item Сохранить идентификатор в \texttt{request.state.request\_id}.
\end{enumerate}

Пример заголовка:

\begin{lstlisting}[caption={Пример заголовка}]
X-Request-ID: 23f3f05c-5a0f-432c-a91f-f8b5011f8840
\end{lstlisting}

Проверьте, что \texttt{request\_id} присутствует:

\begin{itemize}
    \item в теле результата \texttt{/predict};
    \item в логах;
    \item в HTTP response headers;
    \item в записях SQLite.
\end{itemize}

\subsection{Задание 2. Реализуйте промежуточный слой (middleware)}

Добавьте middleware FastAPI, которое должно: получить HTTP-запрос, получить/создать \texttt{request\_id}, записать время начала, выполнить endpoint, рассчитать время обработки, добавить headers, сформировать HTTP-ответ, который содержит \texttt{X-Request-ID}, \texttt{X-Process-Time}.

Проверьте через Swagger: сравните \texttt{X-Process-Time} для \texttt{/health}, \texttt{/predict} и \texttt{/predict-with-llm}.

\subsection{Задание 3. Добавьте обработку ошибок API}

Реализуйте обработку ошибок:

\begin{itemize}
    \item 400 — нарушение бизнес-правил;
    \item 404 — результат не найден;
    \item 409 — неоднозначный результат для batch request\_id;
    \item 422 — ошибка Pydantic;
    \item 429 — очередь переполнена;
    \item 500 — внутренняя ошибка приложения;
    \item 502 — ошибка внешнего LLM-провайдера;
    \item 503 — модель, БД или внешний сервис недоступны;
    \item 504 — timeout.
\end{itemize}

Ответ с ошибкой должен содержать \texttt{request\_id}, если запрос уже прошёл через middleware.

Пример:

\begin{lstlisting}[caption={Пример ответа с ошибкой}]
{
    "detail": "Prediction storage is unavailable",
    "request_id": "..."
}
\end{lstlisting}

\subsection{Задание 4. Подключите SQLite через aiosqlite}

Создайте БД \texttt{agro\_scoring.db}. И таблицу \texttt{predictions} с полями: \texttt{id}, \texttt{request\_id}, \texttt{farm\_id}, \texttt{risk\_score}, \texttt{risk\_level}, \texttt{recommendation}, \texttt{model\_version}, \texttt{source}, \texttt{created\_at}.

\textbf{Ответьте:} почему блокирующий \texttt{sqlite3.execute()} нежелателен внутри асинхронного FastAPI endpoint?

\subsection{Задание 5. Сохраните результат /predict в SQLite}

Модифицируйте \texttt{POST /predict} так, чтобы после расчёта \texttt{risk\_score}, \texttt{risk\_level} и \texttt{recommendation} результат сохранялся в БД.

Клиент получает «200 OK» только после успешного сохранения результата.

\subsection{Задание 6. Реализуйте чтение прогнозов из БД}

\begin{itemize}
    \item \texttt{GET /predictions}
    \item \texttt{GET /predictions/\{request\_id\}}
    \item \texttt{GET /requests/\{request\_id\}/predictions}
\end{itemize}

\texttt{GET /predictions} должен получать историю из SQLite, а не из словаря Python.

\texttt{GET /predictions/\{request\_id\}} должен:

\begin{itemize}
    \item вернуть прогноз при наличии одной записи;
    \item вернуть 404, если запись отсутствует;
    \item вернуть 409, если \texttt{request\_id} связан с batch и содержит несколько predictions.
\end{itemize}

Для batch-запроса используйте \texttt{GET /requests/\{request\_id\}/predictions}.

\subsection{Задание 7. Реализуйте пагинацию результатов}

Создайте endpoint:

\begin{lstlisting}[caption={Пагинация по дню недели}]
GET /predictions/by-weekday
\end{lstlisting}

с параметрами: \texttt{weekday}, \texttt{page}, \texttt{per\_page}.

Пример:

\begin{lstlisting}[caption={Пример запроса}]
GET /predictions/by-weekday?weekday=monday&page=1&per_page=5
\end{lstlisting}

Endpoint должен вернуть первые пять записей понедельника и метаинформацию:

\begin{lstlisting}[caption={Пример ответа пагинации}]
{
    "data": [],
    "meta": {
        "weekday": "monday",
        "current_page": 1,
        "per_page": 5,
        "total_pages": 4,
        "total_items": 18,
        "next_page": "...",
        "prev_page": null
    }
}
\end{lstlisting}

\subsection{Задание 8. Реализуйте явный пакетный инференс (explicit batch inference)}

Создайте:

\begin{lstlisting}[caption={Explicit batch endpoint}]
POST /batch-predict
\end{lstlisting}

Endpoint должен принимать несколько объектов \texttt{FarmRequest}:

\begin{lstlisting}[caption={Тело batch-запроса}]
{
    "items": [
        {"farm_id": "FARM-001", "...": "..."},
        {"farm_id": "FARM-002", "...": "..."},
        {"farm_id": "FARM-003", "...": "..."}
    ]
}
\end{lstlisting}

Ограничьте размер списка:

\begin{lstlisting}[caption={Ограничение размера batch}]
MAX_BATCH_SIZE = 8
\end{lstlisting}

Пакет должен обрабатываться одним вызовом \texttt{model\_predict\_batch(...)}, а результаты сохраняться в SQLite одной транзакцией:

\begin{lstlisting}[caption={Пакетное сохранение}]
await db.executemany(...)
await db.commit()
\end{lstlisting}

\textbf{Ответьте:} почему один \texttt{executemany()} предпочтительнее восьми отдельных \texttt{INSERT + COMMIT}?

\subsection{Задание 9. Реализуйте очередь и динамическую пакетную обработку (dynamic batching)}

Создайте \texttt{asyncio.Queue} и endpoint:

\begin{lstlisting}[caption={Dynamic batching endpoint}]
POST /predict-queued
\end{lstlisting}

Клиент по-прежнему должен отправлять один объект, однако сервер должен помещать его в очередь
Worker должен: получить первый запрос, начать формировать batch, при достижении \texttt{MAX\_BATCH\_SIZE} (или истечении \texttt{MAX\_BATCH\_WAIT\_MS}) вызвать модель.

Используйте:

\begin{lstlisting}[caption={Параметры dynamic batching}]
MAX_BATCH_SIZE = 8
MAX_BATCH_WAIT_MS = 50
\end{lstlisting}

и объект \texttt{Future}, чтобы вернуть соответствующему запросу именно его результат.

\subsection{Задание 10. Реализуйте управление потоком запросов backpressure}

Размер очереди ограничен:

\begin{lstlisting}[caption={Ограничение очереди}]
QUEUE_MAX_SIZE = 100
\end{lstlisting}

Если очередь заполнена, endpoint должен вернуть \textbf{429 Too Many Requests}.

Добавьте:

\begin{lstlisting}[caption={Endpoint /queue-status}]
GET /queue-status
\end{lstlisting}

Endpoint должен возвращать:

\begin{lstlisting}[caption={Ответ /queue-status}]
{
    "queue_size": 4,
    "queue_max_size": 100,
    "max_batch_size": 8,
    "max_batch_wait_ms": 50,
    "inference_timeout_seconds": 2.0
}
\end{lstlisting}

\textbf{Ответьте:} что произойдёт с памятью и latency, если очередь сделать неограниченной?

\subsection{Задание 11. Реализуйте timeout инференса}

Для ограничения времени выполнения модели используйте \texttt{asyncio.wait\_for(...)} и \texttt{INFERENCE\_TIMEOUT\_SECONDS}.

\begin{lstlisting}[caption={Таймаут инференса}]
result = await asyncio.wait_for(
    model_predict(data),
    timeout=INFERENCE_TIMEOUT_SECONDS
)
\end{lstlisting}

При превышении времени вернуть \textbf{504 Gateway Timeout}.

Протестируйте изменение:

\begin{lstlisting}[caption={Тестовое значение таймаута}]
INFERENCE_TIMEOUT_SECONDS = 0.001
\end{lstlisting}

и зафиксируйте результат.

\subsection{Задание 12. Подключите внешнюю LLM}

Подключите Vireonix (без API-ключа и авторизации) через OpenAI-совместимый endpoint:

\begin{lstlisting}[caption={Endpoint LLM}]
POST https://vireonix.ai/v1/chat/completions
\end{lstlisting}

\begin{lstlisting}[caption={Параметры LLM}]
LLM_BASE_URL = "https://vireonix.ai/v1"
LLM_MODEL = "auto"
\end{lstlisting}

\subsection{Задание 13. Реализуйте POST /llm/chat}

Создайте endpoint:

\begin{lstlisting}[caption={Endpoint /llm/chat}]
POST /llm/chat
\end{lstlisting}

Принимать:

\begin{lstlisting}[caption={Тело запроса /llm/chat}]
{
    "prompt": "Объясни, что такое агроскоринг",
    "model": "auto",
    "temperature": 0.3,
    "max_tokens": 300
}
\end{lstlisting}

Ответ должен содержать: \texttt{request\_id}, \texttt{requested\_model}, \texttt{returned\_model}, \texttt{answer}, \texttt{finish\_reason}, \texttt{latency\_ms}.

Проверьте endpoint через Swagger.

\subsection{Задание 14. Реализуйте /predict-with-llm}

Создайте endpoint:

\begin{lstlisting}[caption={Endpoint /predict-with-llm}]
POST /predict-with-llm
\end{lstlisting}

LLM не должна изменять \texttt{risk\_score} и \texttt{risk\_level}. Она должна только сформировать текстовое объяснение результата.

\subsection{Задание 15. Проведите тестирование API}

Заполните следующую таблицу.

\begin{longtable}{|c|p{5.2cm}|c|c|c|}
\caption{Таблица тест-кейсов с ожидаемыми HTTP-кодами}
\label{tab:test_cases}\\
\hline
\textbf{№} & \textbf{Запрос} & \textbf{Ожидаемый HTTP-код} & \textbf{Полученный код} & \textbf{Результат} \\
\hline
\endfirsthead
\hline
\textbf{№} & \textbf{Запрос} & \textbf{Ожидаемый HTTP-код} & \textbf{Полученный код} & \textbf{Результат} \\
\hline
\endhead
1  & GET /health & 200 & & \\
2  & корректный POST /predict & 200 & & \\
3  & area\_ha = -100 & 422 & & \\
4  & неизвестный регион & 400 & & \\
5  & существующий request\_id & 200 & & \\
6  & неизвестный request\_id & 404 & & \\
7  & корректный POST /batch-predict & 200 & & \\
8  & batch больше MAX\_BATCH\_SIZE & 422 & & \\
9  & POST /predict-queued & 200 & & \\
10 & переполненная очередь & 429 & & \\
11 & GET /predictions/by-weekday?page=1\&per\_page=5 & 200 & & \\
12 & page = -1 & 422 & & \\
13 & неизвестный weekday & 400 & & \\
14 & корректный POST /llm/chat & 200 & & \\
15 & недоступный LLM / provider error & 502/503/504 & & \\
16 & корректный POST /predict-with-llm & 200 & & \\
\hline
\end{longtable}

\subsection{Задание 16. Проанализируйте API как контракт}

Заполните таблицу для основных endpoints.

\begin{longtable}{|p{3.2cm}|p{1.4cm}|p{2.4cm}|p{3cm}|p{1.2cm}|p{2.4cm}|}
\caption{Таблица API-контракта}
\label{tab:api_contract}\\
\hline
\textbf{Endpoint} & \textbf{Method} & \textbf{Input} & \textbf{Output} & \textbf{Success} & \textbf{Возможные ошибки} \\
\hline
\endfirsthead
\hline
\textbf{Endpoint} & \textbf{Method} & \textbf{Input} & \textbf{Output} & \textbf{Success} & \textbf{Возможные ошибки} \\
\hline
\endhead
/health & GET & --- & status & 200 & 500 \\
/predict & POST & FarmRequest & PredictionResponse & 200 & 400, 422, 503, 504 \\
/batch-predict & POST & BatchRequest & BatchPredictionResponse & 200 & 400, 422, 503, 504 \\
/predict-queued & POST & FarmRequest & PredictionResponse & 200 & 400, 422, 429, 503, 504 \\
/predictions & GET & query & list & 200 & 400, 422, 503 \\
/predictions/by-weekday & GET & weekday, page, per\_page & paginated response & 200 & 400, 404, 422 \\
/llm/chat & POST & LLMChatRequest & LLMChatResponse & 200 & 422, 502, 503, 504 \\
/predict-with-llm & POST & FarmRequest & PredictWithLLMResponse & 200 & 400, 422, 502, 503, 504 \\
\hline
\end{longtable}

\section{Требования к отчёту}

Отчёт должен содержать:

\begin{enumerate}
    \item Цель работы.
    \item Краткое описание расширенной архитектуры AI-сервиса.
    \item Ссылку на исходный код.
    \item Скриншот Swagger UI.
    \item Скриншот POST /predict.
    \item Скриншот записи результата в SQLite или ответа GET /predictions.
    \item Скриншот POST /batch-predict.
    \item Скриншот GET /queue-status.
    \item Скриншот пагинации по дню недели.
    \item Скриншот POST /llm/chat.
    \item Скриншот POST /predict-with-llm.
    \item Таблицу тест-кейсов.
    \item Таблицу API-контракта.
    \item Краткий вывод о назначении \texttt{request\_id}, batching, backpressure, БД и LLM в составе AI-сервиса.
\end{enumerate}

\section{Критерии оценки}

\begin{longtable}{|p{11cm}|c|}
\caption{Критерии оценки}
\label{tab:criteria}\\
\hline
\textbf{Критерий} & \textbf{Баллы} \\
\hline
\endfirsthead
\hline
\textbf{Критерий} & \textbf{Баллы} \\
\hline
\endhead
Приложение запускается, Swagger/OpenAPI доступны & 5 \\
Реализован middleware с request\_id и X-Process-Time & 10 \\
request\_id проходит через обработку запроса, логирование и БД & 5 \\
SQLite подключена через aiosqlite и await & 10 \\
/predict сохраняет результат в SQLite & 5 \\
Реализовано чтение результатов и поиск по request\_id & 5 \\
Реализована пагинация LIMIT/OFFSET + COUNT(*) & 10 \\
Реализован explicit batching /batch-predict & 10 \\
Реализована очередь asyncio.Queue и dynamic batching & 10 \\
Реализованы backpressure и timeout & 5 \\
Реализован /queue-status и наблюдаемость очереди & 5 \\
Реализовано асинхронное подключение LLM через httpx.AsyncClient & 5 \\
Реализован /llm/chat & 5 \\
Реализован /predict-with-llm с разделением scoring и explanation & 5 \\
Проведено тестирование и заполнена таблица API-контракта & 5 \\
\hline
\textbf{Итого} & \textbf{100} \\
\hline
\end{longtable}

\end{document}